这种失配对协方差自适应后端并非抽象的几何差异，而会落到求解过程本身。以 CMA-ES 为例，若一个子问题同时含有极平坦和极陡峭的方向，步长便无法兼顾两者：能够沿平坦谷底推进的步长对陡峭方向过大，避免陡峭方向发散的步长又会让平坦方向推进极慢。协方差自适应机制可以逐步缓解失配，不过这一过程需要充足采样和学习。在有限预算下，分组造成的局部条件性仍会实质性增加后端求解难度，而现有分组研究对这一实际困难着墨甚少。